\section{Baseline Implementations}\label{app:baselines}

\looseness=-1 Every baseline runs under the same Codex \texttt{gpt-5.5} backbone at high reasoning effort.
The trajectory-only family (Dynamic Cheatsheet, ReasoningBank, and Sleep-time Compute) shares \method{}'s Coreset Selection budget of 10 training tasks and 3 candidates.
Baselines differ only in what the offline phase persists and how the solver consumes it.
This appendix documents each.

\subsection{Dynamic Cheatsheet \citep{suzgun2025dynamiccheatsheet}}

Dynamic Cheatsheet curates a running list of reusable facts and procedures harvested from past trajectories.

\noindent\textbf{Persistence layer.}
A single markdown file lives inside the harness, structured as \texttt{<memory\_item>} blocks containing a description, a worked example, and a usage count.

\noindent\textbf{Offline phase.}
A curator agent iterates over the selected training tasks in order.
For each task it reads the current cheatsheet plus the solve transcript and rewrites the cheatsheet in place.

\noindent\textbf{Online consumption.}
The cheatsheet is part of the harness, and the solver reads it as part of normal harness lookup.
No additional prompt prefix is injected.

\noindent\textbf{Faithfulness.}
Upstream Dynamic Cheatsheet updates the cheatsheet after every task, so the next task's solver sees the update.
Our default configuration runs solves for all selected training tasks before invoking the curator, then commits one cheatsheet.
The single-task configuration that matches the upstream per-example update is available as a sensitivity setting, while Table~\ref{tab:main} uses the default.

\subsection{ReasoningBank \citep{ouyang2026reasoningbank}}

ReasoningBank stores reusable reasoning patterns extracted from past trajectories and retrieves them on demand at inference time.

\noindent\textbf{Persistence layer.}
A JSONL bank of items, each with title, description, body, and a precomputed embedding.
The bank lives outside the harness directory because solve-time access is by similarity retrieval, not by harness materialization.

\noindent\textbf{Offline phase.}
For each selected training task we solve it, ask the judge whether the trajectory was successful, then run an extraction prompt that distills reusable reasoning patterns from the trajectory.
Extracted items are appended to the bank, and embeddings are computed locally.

\noindent\textbf{Online consumption.}
At each held-out solve, we encode the task description with the same encoder, retrieve the top-$n$ items by cosine similarity, render them as a memory preamble prepended to the solve instructions.
Top-$n$ is fixed at 1.

\looseness=-1\noindent\textbf{Faithfulness.}
We replace the upstream Gemini embedding with \texttt{BAAI/bge-large-en-v1.5} \citep{xiao2023cpack}, a 1024-dimensional locally executed sentence encoder, so that every baseline that uses retrieval shares the same encoder, dimensionality, and storage backend.
This removes a confound between memory-bank quality and embedding-provider quality.

\subsection{Sleep-time Compute \citep{lin2025sleeptime}}

Sleep-time Compute preprocesses past traces into compact notes that are prepended to the agent's context at inference time, following the Letta implementation \citep{packer2023memgpt}.

\noindent\textbf{Persistence layer.}
A set of bounded markdown memory blocks stored inside the harness, edited with structured edit tools (insert, replace, rethink, finish) that mirror upstream Letta semantics.

\noindent\textbf{Offline phase.}
A sleep-time agent iterates over the selected training tasks.
For each task it reads the task's prompt and trajectories, then issues a sequence of memory-edit tool calls to update the harness's memory blocks.
The system prompt is the upstream Letta system prompt, used verbatim.

\noindent\textbf{Online consumption.}
The solver reads the memory blocks as part of the harness, the same way it reads any harness file.

\looseness=-1\noindent\textbf{Scope.}
Edits accumulate across tasks, and the sleep-time agent works against a single evolving memory state task after task.
We do not reset memory between training tasks.

\subsection{Meta-Harness \citep{lee2026metaharness}}

Meta-Harness is the validation-feedback reference, where a meta-agent rewrites the harness, and the rewrites are scored against a labeled validation set.

\noindent\textbf{Persistence layer.}
A sequence of candidate harness directories paired with a search history recording per-candidate validation pass rates.

\noindent\textbf{Offline phase.}
An outer loop alternates between a proposer that reads the search history and emits a new harness, and an evaluator that grades the proposed harness on the search-task set using the dataset's ground-truth grader.
Our re-implementation's default budget is 20 outer iterations $\times$ 3 candidates per iteration $\times$ 2 solve trials per task, far above the budget we give \method{}; upstream reports evaluating roughly 60 candidate harnesses over 20 iterations.

\looseness=-1\noindent\textbf{Validation-feedback footprint.}
The proposer reads each candidate's per-task scores, mean score, and pass rate from the history, and the grader signal directly shapes the next proposal.
This is the hallmark of the validation-feedback family and the axis along which Meta-Harness is incomparable to the trajectory-only baselines.

\noindent\textbf{Matched-budget configuration.}
For Table~\ref{tab:meta-harness} we run the outer loop for one iteration with 3 candidates and 1 solve trial per task, so the candidate count matches \method{}'s $N=3$.
The validation-grade cost is what Table~\ref{tab:meta-harness} reports.

\looseness=-1\noindent\textbf{Shared configuration.}
All four baselines and \method{} start from the same empty harness for a given dataset, use the same coreset of selected training tasks, and are graded against the same held-out split with the same grader.
The only deliberate axis of variation is what the offline phase persists and how the solver consumes it.
Meta-Harness additionally consumes validation grades, and this is the axis Table~\ref{tab:meta-harness} measures.
